\begin{textsample}{NEG Dim 6 – global\_south\_2025\_02 – Score 1.00 – global\_south\_2025\_02/120737618.txt}  \label{ex:f6_neg_038}
As ceasefire progresses, Rafah border crossing in Gaza reopens for first time since May Egypt ' s Al-Qahera television broadcast footage of at least two Palestinian Red Cross ambulances arriving at the crossing gate.
Ambulances drive as the Rafah border crossing between Egypt and the Gaza Strip reopens, amid a ceasefire between Israel and Hamas, in Rafah, Egypt, February 1, 2025.
REUTERS/Mohamed Abd El Ghany In another significant breakthrough that bolsters the ceasefire deal Israel and Hamas agreed to earlier this month, Rafah border crossing, the primary entry and exit point for the Palestinian territory, reopened on Saturday for the first time since May 2024.
Israel agreed to reopen the crossing after Hamas released the last living female hostages in Gaza.
Egypt ' s Al-Qahera television broadcast footage of at least two Palestinian Red Cross ambulances arriving at the crossing gate.
Several children were transferred to Egyptian ambulances and taken to hospitals in the nearby city of el-Arish and other locations.
Among them, a young girl who had undergone a @ @ @ @ @ @ @ @ @ @, a report by Associated Press stated.
According to Gaza ' s Health Ministry, around 60 family members accompanied the children.
This marks the start of what is intended to be regular medical evacuations from Gaza for treatment abroad.
Palestinians wounded in the Israeli bombardment of the Gaza Strip wait before crossing the Rafah border into Egypt, as wounded and sick Palestinians are allowed to leave the Gaza Strip for medical treatment, in Khan Younis, Saturday Feb. 1, 2025. ( AP Photo/Jehad Alshrafi ) Israel ' s military campaign against Hamas, launched in response to the October 7, 2023, attack on southern Israel, has severely damaged Gaza ' s healthcare system.
More than 110,000 Palestinians have been wounded, and most hospitals are out of operation, the Palestinian Health Ministry reports.
Why Rafah is Important?
Rafah is Gaza ' s only crossing that does not enter Israel.
Israeli forces closed it in early May after seizing it during an offensive on the southern city.
Egypt subsequently shut its side of @ @ @ @ @ @ @ @ @ @ war, Rafah served as a critical passage for Palestinians seeking medical treatment abroad, as Gaza ' s healthcare system was already under strain due to a 15-year Israeli-Egyptian blockade aimed at containing Hamas.
Reopening the crossing required complex negotiations between Israeli, Egyptian, and Palestinian officials.
Israeli troops remain stationed at the crossing and in the Philadelphia Corridor, the strip of land along the border.
Israel has refused to allow Hamas to regain control of the crossing, accusing the group of using tunnels to smuggle weapons.
Egypt has insisted that it destroyed the tunnels and halted smuggling years ago.
Hundreds of Egyptians gather in front of the Rafah border crossing between Egypt and the Gaza Strip, to show solidarity with the Palestinians and to reject the proposed displacement plans, Friday, Jan. 31, 2025. ( AP Photo/Mohammed Arafat ) Israel has also rejected the prospect of the Palestinian Authority ( PA ) officially running the crossing.
Instead, Palestinians who previously served as PA border officers will staff the crossing but will not @ @ @ @ @ @ @ @ @ @ them to ensure they have no ties to Hamas, according to a European diplomat.
European Union monitors will also be present, as they were before 2007.
Meanwhile, negotiations for the second phase of the ceasefire deal -- focused on a permanent ceasefire, Israel ' s full withdrawal from Gaza, and the release of remaining hostages -- are set to begin on Monday.
However, Israel remains opposed to the PA taking control of postwar Gaza.

% matched lemmas: 
\end{textsample}
